\appendix
\section{Proofs and Supplementary Derivations}

This appendix collects the complete measure-theoretic proofs used by the manuscript. Integrability conditions needed to define each expectation or covariance are stated explicitly at the relevant result rather than being left to a blanket finiteness convention.

\subsection{Covariance identity}

\begin{theorem}[QBS covariance identity]
Let $S\ge0$ with $0<\mathbb E[S]<\infty$. For any random variable $X$ such that
\[
\mathbb E[|X|]<\infty,
\qquad
\mathbb E[|X|S]<\infty,
\]
we have
\[
\mathbb E_{\mathrm{FP}}[X]-\mathbb E[X]
=
\frac{\operatorname{Cov}(X,S)}{\mathbb E[S]}.
\]
\end{theorem}

\begin{proof}
By definition,
\[
\mathbb E_{\mathrm{FP}}[X]
=
\frac{\mathbb E[XS]}{\mathbb E[S]}.
\]
Using
\[
\mathbb E[XS]
=
\mathbb E[X]\mathbb E[S]+\operatorname{Cov}(X,S),
\]
we obtain
\[
\mathbb E_{\mathrm{FP}}[X]
=
\mathbb E[X]
+
\frac{\operatorname{Cov}(X,S)}{\mathbb E[S]}.
\]
Subtracting $\mathbb E[X]$ gives the result.
\end{proof}

Because $\mathbb E[S]>0$, the sign criterion is exact:
\[
\mathbb E_{\mathrm{FP}}[X]>\mathbb E[X]
\quad\Longleftrightarrow\quad
\operatorname{Cov}(X,S)>0.
\]
Equality holds exactly when $\operatorname{Cov}(X,S)=0$.

\subsection{Tail probability identity}

\begin{theorem}[Tail identity]
For every threshold $c$,
\[
P_{\mathrm{FP}}(X\ge c)-P(X\ge c)
=
\frac{\operatorname{Cov}(\mathbf 1_{\{X\ge c\}},S)}{\mathbb E[S]}.
\]
\end{theorem}

\begin{proof}
Apply the covariance identity to the bounded random variable
\[
Z=\mathbf 1_{\{X\ge c\}}.
\]
Then $\mathbb E[Z]=P(X\ge c)$ and $\mathbb E_{\mathrm{FP}}[Z]=P_{\mathrm{FP}}(X\ge c)$. Because $Z$ is bounded, $\mathbb E[|Z|]<\infty$ and $\mathbb E[|Z|S]\le\mathbb E[S]<\infty$.
\end{proof}

\subsection{Monotone accessibility and first-order stochastic dominance}

\begin{theorem}[Monotone accessibility implies FOSD]
Let
\[
g(x)=\mathbb E[S\mid X=x]
\]
be a version of the conditional expectation that is nondecreasing on the support of $X$. Then
\[
F_{\mathrm{FP}}(c)\le F(c)
\]
for every real $c$.
\end{theorem}

\begin{proof}
Fix $c$ and define
\[
f_c(x)=\mathbf 1_{\{x\le c\}}.
\]
The covariance identity gives
\[
F_{\mathrm{FP}}(c)-F(c)
=
\frac{\operatorname{Cov}(f_c(X),S)}{\mathbb E[S]}.
\]
Since $f_c(X)$ is measurable with respect to $X$,
\[
\operatorname{Cov}(f_c(X),S)
=
\operatorname{Cov}(f_c(X),\mathbb E[S\mid X])
=
\operatorname{Cov}(f_c(X),g(X)).
\]
Let $X'$ be an independent copy of $X$. For integrable functions $a$ and $b$,
\[
2\operatorname{Cov}(a(X),b(X))
=
\mathbb E[(a(X)-a(X'))(b(X)-b(X'))].
\]
Here $f_c$ is nonincreasing and $g$ is nondecreasing. Therefore
\[
(f_c(X)-f_c(X'))(g(X)-g(X'))\le 0
\]
almost surely. Hence
\[
\operatorname{Cov}(f_c(X),g(X))\le 0.
\]
Division by the positive quantity $\mathbb E[S]$ proves
\[
F_{\mathrm{FP}}(c)\le F(c).
\]
\end{proof}

\paragraph{Equality condition.}
At a fixed threshold $c$,
\[
F_{\mathrm{FP}}(c)=F(c)
\]
if and only if
\[
\operatorname{Cov}(\mathbf 1_{\{X\le c\}},g(X))=0.
\]
A sufficient condition for equality at every threshold is that $g(X)$ is almost surely constant.

\paragraph{Strictness condition.}
At threshold $c$, strict dominance holds whenever two independent copies $X,X'$ have positive probability of satisfying
\[
X\le c<X'
\]
and simultaneously
\[
g(X)<g(X').
\]
On this event the covariance-pair product is strictly negative, and therefore
\[
F_{\mathrm{FP}}(c)<F(c).
\]

\paragraph{Boundary of the theorem.}
Positive Pearson correlation between $X$ and $S$ does not imply FOSD. The proof requires the stronger ordering condition that $g(x)$ be nondecreasing. Nonmonotone accessibility may generate CDF crossings even when the first-person mean is larger.

\subsection{Recognition decomposition}

Let recognition state $R\in\{0,1\}$ induce policy-dependent pairs $(U_R,S_R)$. For each $R$, assume
\[
0<\mathbb E[S_R]<\infty,
\qquad
\mathbb E[|U_R|]<\infty,
\qquad
\mathbb E[|U_R|S_R]<\infty.
\]
Define
\[
V_R
=
\frac{\mathbb E[U_RS_R]}{\mathbb E[S_R]},
\qquad
Q(U,S)
=
\frac{\operatorname{Cov}(U,S)}{\mathbb E[S]}.
\]

\begin{theorem}[Recognition decomposition]
\[
V_1-V_0
=
\mathbb E[U_1-U_0]
+
Q(U_1,S_1)-Q(U_0,S_0).
\]
\end{theorem}

\begin{proof}
Apply the covariance identity separately to each recognition state:
\[
V_R
=
\mathbb E[U_R]+Q(U_R,S_R).
\]
Subtract the expression for $R=0$ from that for $R=1$.
\end{proof}

If the pre-recognition accessibility is constant,
\[
S_0\equiv 1,
\]
then $Q(U_0,S_0)=0$ and
\[
V_1-V_0
=
\mathbb E[U_1-U_0]
+
\frac{\operatorname{Cov}(U_1,S_1)}{\mathbb E[S_1]}.
\]
The first term is the ordinary trajectory/policy effect; the second is the first-person conditioning contribution.

\paragraph{Recognition-label null.}
If recognition changes neither trajectory nor accessibility,
\[
U_1=U_0,
\qquad
S_1=S_0
\]
almost surely, then $V_1-V_0=0$ exactly.

\subsection{Policy--QBS interaction decomposition}

For the general selector-changing identity, additionally assume
\[
\mathbb E[|U_1|S_0]<\infty,
\]
so that the intermediate term $Q(U_1,S_0)$ is defined. Let
\[
D=U_1-U_0,
\qquad
I=Q(U_1,S_1)-Q(U_0,S_0).
\]

\begin{theorem}[Interaction decomposition]
\[
I
=
\frac{\operatorname{Cov}(D,S_0)}{\mathbb E[S_0]}
+
\left[Q(U_1,S_1)-Q(U_1,S_0)\right].
\]
\end{theorem}

\begin{proof}
Add and subtract $Q(U_1,S_0)$:
\[
I
=
\left[Q(U_1,S_0)-Q(U_0,S_0)\right]
+
\left[Q(U_1,S_1)-Q(U_1,S_0)\right].
\]
Since $U_1=U_0+D$ and covariance is linear in its first argument,
\[
Q(U_1,S_0)-Q(U_0,S_0)
=
\frac{\operatorname{Cov}(D,S_0)}{\mathbb E[S_0]}.
\]
Substitution proves the identity.
\end{proof}

For a fixed selector $S_1=S_0=S$,
\[
I
=
\frac{\operatorname{Cov}(D,S)}{\mathbb E[S]}.
\]
Therefore the sign of the interaction is exactly the sign of $\operatorname{Cov}(D,S)$.

\subsection{Adaptive-rescue sign condition}

\begin{corollary}[Adaptive rescue]
Let $B$ denote branch badness. Suppose policy improvement is $D=d(B)$ with $d$ nondecreasing and accessibility is $S=s(B)$ with $s$ nonincreasing. Then
\[
\operatorname{Cov}(D,S)\le 0.
\]
Under a fixed selector, $I\le 0$.
\end{corollary}

\begin{proof}
Let $B'$ be an independent copy of $B$. Then
\[
2\operatorname{Cov}(d(B),s(B))
=
\mathbb E[(d(B)-d(B'))(s(B)-s(B'))].
\]
Because $d$ and $s$ have opposite monotonicity, the integrand is nonpositive almost surely. Hence the covariance is nonpositive.
\end{proof}

Strict negativity holds when the two functions vary nontrivially in opposite directions on a set that can be separated by two independent draws with positive probability.

\subsection{Option value and support preservation}

\begin{proposition}[Costless recognition option value]
If $\Pi_0\subseteq\Pi_1$ and the same value functional $V$ evaluates both policy sets, then
\[
\sup_{\pi\in\Pi_1}V(\pi)
\ge
\sup_{\pi\in\Pi_0}V(\pi).
\]
\end{proposition}

\begin{proof}
The supremum over a superset cannot be smaller than the supremum over a subset.
\end{proof}

\begin{proposition}[Pure reweighting cannot create support]
For fixed policy $\pi$, the first-person measure is absolutely continuous with respect to the base measure:
\[
\mu^{\mathrm{FP}}_\pi\ll\mu.
\]
Thus $\mu(A)=0$ implies $\mu^{\mathrm{FP}}_\pi(A)=0$.
\end{proposition}

\begin{proof}
If $\mu(A)=0$, then $\mathbf 1_A=0$ almost surely. Hence
\[
\mathbb E[\mathbf 1_A S_\pi]=0,
\]
and therefore $\mu^{\mathrm{FP}}_\pi(A)=0$ whenever $\mathbb E[S_\pi]>0$.
\end{proof}

Recognition may nevertheless change the policy and therefore the trajectory map itself. Support preservation applies to pure reweighting under a fixed policy, not to causal policy changes.

\subsection{Zero-accessible-measure boundary}

The normalized first-person measure requires
\[
\mathbb E[S]>0.
\]
If $\mathbb E[S]=0$, the normalized measure is undefined. This is not merely a low-value state; normalization itself fails.

\subsection{Repeated-filter sensitivity identity}

Let $N_B$ count adverse triggers and define
\[
S=\lambda^{N_B},
\qquad
V(\lambda)
=
\frac{\mathbb E[U\lambda^{N_B}]}{\mathbb E[\lambda^{N_B}]}.
\]
Assume differentiation under the expectation is valid. Define the $\lambda$-weighted expectation by
\[
\mathbb E_\lambda[X]
=
\frac{\mathbb E[X\lambda^{N_B}]}{\mathbb E[\lambda^{N_B}]}.
\]
Then
\[
\frac{dV}{d\log\lambda}
=
\operatorname{Cov}_\lambda(U,N_B).
\]
Indeed, writing $N(\lambda)=\mathbb E[U\lambda^{N_B}]$ and $D(\lambda)=\mathbb E[\lambda^{N_B}]$, the quotient rule gives
\[
\lambda\frac{dV}{d\lambda}
=
\frac{\mathbb E[UN_B\lambda^{N_B}]}{D(\lambda)}
-
\frac{N(\lambda)}{D(\lambda)}
\frac{\mathbb E[N_B\lambda^{N_B}]}{D(\lambda)},
\]
which is the weighted covariance.

\subsection{Gaussian minimal model}

Let $(L,Y)$ be jointly standard normal with correlation $\rho$. Let
\[
S
=
\begin{cases}
1, & Y\ge 0,\\
\lambda, & Y<0.
\end{cases}
\]
Because
\[
\mathbb E[L\mid Y]=\rho Y,
\]
we have
\[
\mathbb E[LS]
=
\rho\,\mathbb E[YS].
\]
Using
\[
\mathbb E[Y\mathbf 1_{\{Y\ge0\}}]
=
\frac{1}{\sqrt{2\pi}},
\qquad
\mathbb E[Y\mathbf 1_{\{Y<0\}}]
=
-\frac{1}{\sqrt{2\pi}},
\]
we obtain
\[
\mathbb E[YS]
=
\frac{1-\lambda}{\sqrt{2\pi}}.
\]
Also
\[
\mathbb E[S]
=
\frac{1+\lambda}{2}.
\]
Therefore
\[
\mathbb E_{\mathrm{FP}}[L]
=
\frac{2(1-\lambda)\rho}{(1+\lambda)\sqrt{2\pi}}.
\]
Under the toy parameterization $\lambda=1-q(1-\alpha)$,
\[
\mathbb E_{\mathrm{FP}}[L]
=
\frac{2q(1-\alpha)\rho}{[2-q(1-\alpha)]\sqrt{2\pi}}.
\]
This is an exact result for the Gaussian toy model, not a physical Everett prediction.

\subsection{Separate normalization for multiple observers}

For observer $i$, define
\[
d\mu_i^{\mathrm{FP}}
=
\frac{S_i}{\mathbb E[S_i]}
\,d\mu.
\]
Each observer's first-person measure is normalized separately. These normalized measures are not shares of a single probability budget and therefore need not sum over observer index $i$. This observation is purely measure-theoretic and does not imply that mutually incompatible shared-world events are jointly likely.

\subsection{Reproducibility conventions}

The principal experiments use explicit random seeds and committed source data. GitHub Actions reruns E1--E5, regenerates the publication SVG figures, validates the experiment manifest, and independently compiles the manuscript PDF. Historical exploratory parameter sweeps remain secondary evidence and are not promoted to theorem status.
